% \documentclass[twocolumn, 10pt]{article}
\documentclass[twocolumn,9pt]{extarticle}

\usepackage{geometry}
\geometry{
	a4paper,
	total={6.85in, 9.92in},
	left=0.71in,
	top=0.63in,
}

% \usepackage{hyperref}
% \usepackage[hidelinks]{hyperref}
% \hypersetup{
%     colorlinks = true,   % Use colored text instead of boxes
%     linkcolor = gray,    % Color for internal links (e.g., \ref, \cite)
%     citecolor = blue,    % Color for citations
%     filecolor = magenta, % Color for file links
%     urlcolor = gray      % Color for URLs
% }

\usepackage{graphicx}
\usepackage{xcolor}
\usepackage{caption}

% pseudocode
\usepackage{amsmath}
\usepackage{amssymb}
\usepackage{algorithm}
\usepackage{algpseudocode}
\usepackage{amsmath}

% % utf
% \usepackage[T1]{fontenc}
% \usepackage[utf8]{inputenc}
% \usepackage[polish]{babel}

\usepackage{enumitem}
% \setlist{nosep}
\setlist[itemize]{itemsep=1pt, topsep=1pt, leftmargin=1.2em}

\captionsetup{font={scriptsize}}
\renewcommand{\footnotesize}{\scriptsize}
\pagenumbering{gobble}

\title{Semantic Kademlia?}
\author{Artur Augustyniak}

\makeatletter
\newcommand{\fsize}{\f@size pt }
\newcommand{\textFontName}{\f@family}

\renewcommand{\maketitle}{
\begin{flushleft}
{\noindent\Huge\bf\@title}\break
\end{flushleft}
}

% \renewcommand{\maketitle}{
% \begin{flushleft}
% {\noindent\Huge\bfseries\@title}
% \vspace{1em} % vertical gap
% \end{flushleft}
% }

\makeatother
\usepackage{xurl}                    % better URL breaking
\urlstyle{tt}

% ---- hyperref ONCE, near the end
\usepackage[colorlinks=true,
            linkcolor=gray,
            citecolor=blue,
            filecolor=magenta,
            urlcolor=gray]{hyperref}

% ---- Footnotes: smaller + better layout
\usepackage[hang,splitrule]{footmisc} % hanging number; nice split across columns
\setlength{\footnotemargin}{1.2em}    % space between number and text
\renewcommand\footnotelayout{\fontsize{8}{9.5}\selectfont} % smaller footnote text
\setlength{\skip\footins}{6pt}        % space above footnote area
\setlength{\footnotesep}{0.35\baselineskip} % space between footnotes

\begin{document}
\maketitle

While experimenting with a Kademlia\footnote{\url{https://pdos.csail.mit.edu/~petar/papers/maymounkov-kademlia-lncs.pdf}} simulation, I started wondering: 
since the identifier space is discrete, routing relies on the XOR metric, what would happen if I replace those with vectors (embeddings), and XOR distance with cosine distance?

It turns out that this leads to a working simulation of a network that:
\begin{itemize}
  \item stores data in a decentralized manner,
  \item allows semantic queries (e.g., in natural language),
  \item replicates data near its “position” in embedding space,
  \item adapts routing locally in response to queries and topology changes,
  \item operates in a best-effort mode without guarantees of optimal routing.
\end{itemize}

For brevity:
\begin{itemize}
\item simulation driver: \url{https://github.com/artur-augustyniak/semantic_kademlia/blob/main/main.py},
\item example run log: \url{https://github.com/artur-augustyniak/semantic_kademlia/blob/main/data/network.log}.
\end{itemize}

\vspace{-9pt}
\subsection*{Embeddings without heavylifting}
% \vspace{-6pt}

I did not want to use heavy language models in this toy simulation. 
Instead\footnote{\url{https://github.com/artur-augustyniak/semantic_kademlia/blob/afa05ea7ea70d76b0a6c4db4cc326a81c3acc764/tooling.py\#L65}}, 
I ended up using Feature Hashing\footnote{\url{https://www.usna.edu/Users/cs/crabbe/SD311/2024/lectures/maps/feature.html}}/CountSketch\footnote{\url{https://proceedings.mlr.press/v139/larsen21a/larsen21a.pdf}}vectorizer.

\algrenewcommand\algorithmicindent{0.6em}
\algrenewcommand\algorithmicrequire{\textbf{In:}}
\algrenewcommand\algorithmicensure{\textbf{Out:}}

\begin{algorithm}
% \scriptsize
\captionsetup{font=scriptsize}
\caption{Lexical embedding via CountSketch over character n-grams}
\label{alg:countsketch-embedding}
\begin{algorithmic}
\Require text $t$, embedding dimension $d$
\Ensure vector $\mathbf{v} \in \mathbb{R}^d$

\State $t \gets \text{lowercase}(\text{trim}(t))$
\If{$t$ is empty}
    \State \Return $\mathbf{0} \in \mathbb{R}^d$
\EndIf

\State $t \gets \langle t \rangle$ \Comment{add boundary markers}
\State $\mathbf{v} \gets \mathbf{0} \in \mathbb{R}^d$

\For{$n \in \{2,3,4\}$}
    \For{$i \gets 0$ to $|t| - n$}
        \State $g \gets t[i : i+n]$
        \State $h \gets \text{Hash}(g)$
        \State $j \gets h \bmod d$
        \State $s \gets
        \begin{cases}
            +1 & \text{if } \mathrm{bit}_{31}(h) = 0 \\
            -1 & \text{otherwise}
        \end{cases}$
        \State $v_j \gets v_j + s$
    \EndFor
\EndFor

\State $\mathbf{v} \gets \mathbf{v} / \|\mathbf{v}\|_2$
\State \Return $\mathbf{v}$
\end{algorithmic}
\end{algorithm}

The result is not "semantic" in the sense embedding models\footnote{\url{https://huggingface.co/sentence-transformers/all-MiniLM-L6-v2}} - it a clever approximation of lexical similarity.
Texts with similar structure end up close in the vector space, although there is no deeper representation of meaning.

\vspace{-9pt}
\subsection*{What is actually happening here}
% \vspace{-6pt}

The XOR metric has strong properties - it imposes a hierarchical structure and enables predictable narrowing of the search space. 
In embedding space, these properties disappear: the fact that vector $B$ is closer to $C$ than $A$ does not imply that its neighbors will also be closer with respect to $C$. 

So, I do not expect stable, globally correct routing. I treat the network as an adaptive system that builds its structure through queries and data replication\footnote{\url{https://github.com/artur-augustyniak/semantic_kademlia/blob/afa05ea7ea70d76b0a6c4db4cc326a81c3acc764/node.py\#L114}}.
In local neighborhoods, a form of proximity transitivity often holds, which allows the network to gradually improve answer quality over time\footnote{\url{https://github.com/artur-augustyniak/semantic_kademlia/blob/afa05ea7ea70d76b0a6c4db4cc326a81c3acc764/node.py\#L268}}.

The cost is longer search paths: queries require higher TTL values and more often end up in dead ends. However, I this effect is amortized over time - the routing table gradually move itself towards "hot" regions.

This also motivates one nasty trick: pruning the routing table with respect to the most recently observed query or key\footnote{\url{https://github.com/artur-augustyniak/semantic_kademlia/blob/afa05ea7ea70d76b0a6c4db4cc326a81c3acc764/node.py\#L69}}. 
This helps counteract the drift caused by random exploration, which is necessary to avoid blindness to new 
nodes\footnote{\url{https://github.com/artur-augustyniak/semantic_kademlia/blob/afa05ea7ea70d76b0a6c4db4cc326a81c3acc764/node.py\#L418}}, but at the same time degrades locality.

These issues become more pronounced as dimensionality increases. The curse of dimensionality\footnote{\url{https://en.wikipedia.org/wiki/Curse_of_dimensionality}} causes distances to become less discriminative, weakening 
the routing signal. At the same time, higher dimensionality reduces collisions in the embedding, improving representation quality\footnote{\url{https://github.com/artur-augustyniak/semantic_kademlia/blob/afa05ea7ea70d76b0a6c4db4cc326a81c3acc764/main.py\#L26}}, 
but worsening geometric properties relevant for routing - leading to a trade-off between embedding quality and the “navigability” of the space.

Additionally, unlike classical Kademlia where hashing enforces a near-uniform distribution of data, here semantic clustering emerges naturally, leading to uneven load.
At the same time, this makes the system manipulable in an interesting way: since node identifiers are vectors, one can deliberately choose an ID that lies close to a particular region of the embedding space, effectively making a node “attract” specific kinds of data.

I came across prior work on similar ideas\footnote{\url{https://gist.github.com/ostafen/a556180db7d4c41abb325b5ae5e13ca4}}\footnote{\url{https://arxiv.org/html/2502.10151v1}}.
These are significantly more formal and comprehensive than the simple simulation presented here, and they focus on addressing locality and clustering issues.

It may be beneficial to impose additional structure - for example, 
via a BSP-like partitioning - in order to introduce a partial hierarchy and improve routing properties.
But Maybe, in some cases an alternative is to accept these properties and explore how they can be turned into a feature rather than a limitation.









\end{document}

